\centerline{\bf \Huge Банки. Жесть. КР}
\centerline{\bf \Huge }

\begin{flushright}
        \scriptsize{\textbf{ww}}
\end{flushright}

\problem{\textbf{1}} В июле 2026 года планируется взять кредит на три года на сумму 800 тыс. рублей. Условия его возврата таковы:

— каждый январь долг будет возрастать на 20\% по сравнению с концом предыдущего года;

— с февраля по июнь каждого года необходимо выплатить одним платежом часть долга;

— платежи в 2027 и в 2028 годах должны быть равны;

— к июлю 2029 года долг должен быть выплачен полностью.

Известно, что общая сумма платежей будет равна 1190,4 тыс. рублей. Чему будет равна
третья выплата?

\problem{\textbf{2}} В июле 2026 года планируется взять кредит в банке на пять лет в размере $S$ тыс рублей.
Условия его возврата таковы:

— каждый январь долг возрастает на 20\% по сравнению с концом предыдущего года;

— с февраля по июнь каждого года необходимо выплатить часть долга;

— в июле 2027, 2028 и 2029 годов долг остаётся равным $S$ тысяч рублей;

— выплаты в 2030 и 2031 годах равны по 360 тысяч рублей;

— к июлю 2031 года долг будет выплачен полностью.

Найдите общую сумму выплат за пять лет.

.

.

.

.

\problem{\textbf{3}} 15 января некоторого года Антон взял кредит на 3 млн рублей на 6 месяцев. Условия его возврата таковы:

– 1-го числа каждого месяца долг возрастает на $r$\% по сравнению с концом предыдущего
месяца;

– со 2-го по 14-е число каждого месяца необходимо выплатить часть долга;

– 15-го февраля, апреля и июня долг должен быть на одну девятую часть от исходной суммы
долга меньше, чем величина долга 15-го числа предыдущего месяца;

– 15-го марта, мая и июля долг должен быть на две девятых части от исходной суммы долга
меньше, чем величина долга 15-го числа предыдущего месяца.

Известно, что общая сумма выплат после полного погашения кредита на 220 тыс. рублей
больше суммы, взятой в кредит. Найдите $r$

\problem{\textbf{4}} В июле 2010-го года планируется взять кредит в банке на сумму 670 тыс. рублей на 11 лет.
Условия его возврата таковы:

– каждый январь долг возрастает на $r$\% по сравнению с концом предыдущего года;

– с февраля по июнь каждого года необходимо выплачивать часть долга;

– в июле каждого года с 2011 по 2019 долг должен быть на 50 тыс. рублей меньше долга на
июль предыдущего года;

– в последние два года должны быть внесены равные платежи;

Чему будет равна общая сумма выплат, если известно, что пятый по счету платеж должен
быть равен 144 тыс. рублей?
